\section{twisted cubic と外部直線：七つの semisimple 節点型}

$T=T_{0123},L=L_{04}$ とする。本節は、節点の固有値・消失次数の幾何学的帰着を入力とし、
そこに残る七つの全 Jacobi 条件を厳密に処理する。
semisimple residue を $R_T=t\partial_t,R_L=\rho t\partial_t$ と標準化する。
従って第21版のモデルでは
\[
z_{36}=1,\ z_{37}=0,\ z_{38}=\rho,\ z_{39}=0.
\]
節点で $B_T=0$ は $z_{30}=z_{31}=z_{25}=0$ である。

両方の conormal bracket が節点で零なら、さらに
\[
z_{32}=0,\quad z_{22}=z_{26},\quad
z_{14}=-z_{28}-z_{35}-(2-\rho)/5.
\]
直線の Hessian 消失次数は二であり、三次 flag の消失次数は $0,1,2$ のいずれか。
節点のもう一つの接線固有値が $1$ なので、残る比は
\[
\rho\in\{-1,-3,3,-3/2,3/2\}.
\]
直線側の bracket が非零の場合の二つの候補は、節点の固有値計算により
$(\rho,k_T)=(-3/2,2),(-2,3)$ である。
$k_T\ge2$ は追加の必要条件
\[
z_{33}=z_{32}=z_{26}=z_{21}=0,\qquad z_{27}=2z_{17}
\]
を与え、直線側 rank-one 条件から $z_{22}=0$。
$k_T\ge3$ ならさらに
\[
z_{17}=z_{27}=z_{18}=z_{34}=z_{23}=0,
\qquad 3z_{28}+2z_{35}=3\rho/5.
\]
これらの各スライス上で、全 Jacobi 係数と全 E8 必要係数を保持した。
変数での除算は行わず、有理数線形結合と一次式の消去だけを用いた。

\begin{center}
\begin{tabular}{ccl}
節点型 & $\rho$ & 結果\\
両 bracket 零 & $-1$ & $z_9z_{21}=1/9$ の一族\\
両 bracket 零 & $-3$ & 二段で単位イデアル\\
両 bracket 零 & $3$ & 二段で単位イデアル\\
両 bracket 零 & $-3/2$ & 三段で単位イデアル\\
両 bracket 零 & $3/2$ & 二段で単位イデアル\\
$k_T\ge2$ & $-3/2$ & 二段で単位イデアル\\
$k_T\ge3$ & $-2$ & 一段で単位イデアル
\end{tabular}
\end{center}

唯一残る一族でも、五段の一次消去の後の全イデアルは正確に
$(z_9z_{21}-1/9)$ である。$a=z_9\ne0,z_{21}=1/(9a)$ と置くと、
退化五次式は
\[
F=-\frac1{90a}
\bigl(3a(x_2^2-x_1x_3)+x_0x_4\bigr)
\bigl(3ax_2(x_2^2-x_1x_3)+x_1^2x_4\bigr).
\]
二つの因子は非零で、それぞれ次数二、三である。
従ってその共通零点の曲面は $\operatorname{Sing}D$ に含まれ、実際の退化因子は正規でない。
この一族は Jacobi 解としては実在する。正規性の条件によって排除する点を区別する。

全証明書は有理数係数の線形結合を各段の入力式と共に保存した。
Python 標準ライブラリのみの独立実装により、元の $P$ から両成分の全局所行列と
全 E8 係数を再構成し、七つの証明書、残存イデアル、上の因数分解を確認した。
一族については直接の全十 Jacobi 括弧と $dF=b(P)$ も検証した。
全1243項目が通り、数値代入や CAS の使用はない。
